\chapter{Coronary Artery Disease}

\section*{Definition and Overview}
Coronary artery disease (CAD) is the narrowing or blockage of the coronary arteries, the blood vessels that supply oxygen-rich blood to the heart muscle. It is usually caused by atherosclerosis, the slow buildup of fatty plaque within the artery walls. CAD can be silent for years and then cause chest pain (angina), shortness of breath, or a heart attack. It is the most common form of heart disease and a leading cause of death worldwide, but it is also highly preventable and treatable with lifestyle measures, medicines, and, when needed, procedures to restore blood flow.

\section*{Epidemiology}
Coronary artery disease is the single largest cause of death globally, and it affects a substantial proportion of the adult population in Australia. It is more common with increasing age and in men, although the gap narrows after menopause in women. Major risk factors include smoking, high blood pressure, high cholesterol, diabetes, obesity, physical inactivity, a family history of early heart disease, and unhealthy diet. Around one in ten Australian adults report having been told they have heart disease, and heart attacks remain a major cause of hospitalisation. However, death rates have fallen substantially over recent decades thanks to better prevention and treatment.

\section*{Aetiology and Pathophysiology}
The central process is atherosclerosis, in which the walls of the coronary arteries thicken and narrow over many years.

\begin{itemize}
    \item \textbf{Endothelial injury:} high blood pressure, smoking, and high blood sugar damage the artery lining
    \item \textbf{Plaque formation:} cholesterol and inflammatory cells accumulate, forming a fatty plaque
    \item \textbf{Plaque growth:} plaques slowly enlarge, narrowing the artery and limiting blood flow
    \item \textbf{Plaque rupture:} a plaque can tear, exposing its contents and triggering a blood clot
    \item \textbf{Heart attack:} a clot that completely blocks a coronary artery cuts off blood supply to part of the heart muscle
    \item \textbf{Ischaemia:} a mismatch between oxygen supply and demand, causing angina and reduced pump function
    \item \textbf{Risk factors:} smoking, hypertension, dyslipidaemia, diabetes, obesity, inactivity, and family history
\end{itemize}

\section*{Clinical Features}
CAD can be asymptomatic, or it can cause symptoms when the heart does not get enough blood.

\begin{itemize}
    \item Chest pain, pressure, squeezing, or heaviness (angina), often brought on by exertion and relieved by rest
    \item Discomfort spreading to the arm, jaw, neck, or back
    \item Shortness of breath, especially with exertion
    \item Fatigue and reduced exercise tolerance
    \item Palpitations or an irregular heartbeat
    \item Nausea, sweating, or dizziness, particularly during a heart attack
    \item Some people, especially women and people with diabetes, have atypical symptoms or none at all
\end{itemize}

\section*{Differential Diagnosis}
Chest pain and breathlessness have many causes that should be considered.

\begin{itemize}
    \item Gastro-oesophageal reflux and oesophagitis
    \item Musculoskeletal chest pain, such as costochondritis or muscle strain
    \item Anxiety, panic attacks, and hyperventilation
    \item Pulmonary embolism, pneumonia, or pleurisy
    \item Aortic stenosis and other valve diseases
    \item Pericarditis or myocarditis
    \item Aortic dissection, a medical emergency
\end{itemize}

\section*{When to Seek Medical Care}
See a doctor promptly if you experience chest pain with exertion, unexplained breathlessness, or other possible symptoms of heart disease, especially if you have risk factors. An ECG and blood tests can help establish the cause.

\textbf{Call 000 (or local emergency services) for:}
\begin{itemize}
    \item Chest pain at rest, or pain lasting longer than a few minutes
    \item Pain spreading to the arm, jaw, neck, or back
    \item Shortness of breath, sweating, nausea, or feeling faint with chest pain
    \item Sudden weakness, numbness, or difficulty speaking (possible stroke)
    \item Collapse or loss of consciousness
\end{itemize}

\section*{Diagnosis and Investigations}
The diagnosis is based on symptoms, risk factors, and the results of tests.

\begin{itemize}
    \item \textbf{History and examination:} description of symptoms, risk factors, and family history
    \item \textbf{ECG:} records the heart's electrical activity at rest and may show prior or current damage
    \item \textbf{Blood tests:} cholesterol, blood sugar, and, in acute situations, cardiac troponin
    \item \textbf{Exercise stress test:} ECG and blood pressure monitored during exercise
    \item \textbf{Echocardiography:} ultrasound to assess heart structure and pump function
    \item \textbf{Coronary calcium score:} a CT scan quantifying calcium in the coronary arteries
    \item \textbf{CT coronary angiography or invasive coronary angiography:} direct imaging of the coronary arteries
\end{itemize}

\section*{Management}
Management aims to relieve symptoms, slow the disease, and prevent heart attacks.

\begin{itemize}
    \item \textbf{Lifestyle:} smoking cessation, heart-healthy diet, regular exercise, and weight management
    \item \textbf{Antiplatelet therapy:} low-dose aspirin or other agents to reduce clotting
    \item \textbf{Statins:} cholesterol-lowering medicines that stabilise plaques
    \item \textbf{Blood pressure control:} antihypertensive medicines to protect the arteries
    \item \textbf{Anti-anginal medicines:} beta-blockers, calcium channel blockers, or nitrates
    \item \textbf{Diabetes management:} tight control of blood sugar
    \item \textbf{Percutaneous coronary intervention:} angioplasty and stenting for significant narrowings
    \item \textbf{Coronary artery bypass grafting:} surgery for extensive or complex disease
    \item \textbf{Cardiac rehabilitation:} supervised exercise, education, and support
\end{itemize}

\section*{Complications}
\begin{itemize}
    \item Heart attack (myocardial infarction), with damage to heart muscle
    \item Heart failure, when the pump function of the heart is impaired
    \item Abnormal heart rhythms, including life-threatening arrhythmias
    \item Sudden cardiac death
    \item Angina that limits daily activity and quality of life
    \item Complications of treatment, including bleeding or complications of procedures
    \item Psychological effects, including anxiety and depression after a cardiac event
\end{itemize}

\section*{Prognosis and Outcomes}
The outlook for coronary artery disease has improved dramatically. With modern treatment, most people live for many years after the diagnosis, and death rates from heart attacks continue to fall. Prognosis depends on the extent of disease, pump function, and how well risk factors are controlled. Adherence to medicines and lifestyle change is the single most important factor in long-term outcomes. People who have had a heart attack or have reduced pump function are at higher risk and benefit most from structured secondary prevention.

\section*{Prevention and Self-Care}
\begin{itemize}
    \item Do not smoke and avoid second-hand smoke
    \item Eat a heart-healthy diet rich in vegetables, fruit, wholegrains, and healthy fats
    \item Be physically active for at least 150 minutes of moderate activity each week
    \item Achieve and maintain a healthy weight
    \item Keep blood pressure, cholesterol, and blood sugar in target ranges
    \item Limit alcohol
    \item Take prescribed medicines every day, even when feeling well
    \item Learn the signs of a heart attack and how to respond
    \item Attend regular health checks and cardiac rehabilitation after an event
\end{itemize}

\section*{Sources and Further Reading}
The following reputable sources provide further information for healthcare students and patients seeking reliable, current advice about coronary artery disease.

\begin{itemize}
    \item Heart Foundation Australia. \textit{Coronary heart disease.} https://www.heartfoundation.org.au/
    \item NHS. \textit{Coronary heart disease.} https://www.nhs.uk/conditions/coronary-heart-disease/
    \item Mayo Clinic. \textit{Coronary artery disease: overview.} https://www.mayoclinic.org/diseases-conditions/coronary-artery-disease/
    \item MedlinePlus. \textit{Coronary artery disease health information.} https://medlineplus.gov/coronaryarterydisease.html
    \item World Health Organization. \textit{Cardiovascular diseases fact sheet.} https://www.who.int/
    \item MSD Manual. \textit{Coronary artery disease: professional reference.} https://www.msdmanuals.com/
\end{itemize}
